\documentclass[11pt,a4paper]{article}

% Required packages
\usepackage{graphicx}
\usepackage{float}
\usepackage{caption}
\usepackage{adjustbox}
\usepackage[margin=1.5cm]{geometry}
\usepackage{pdflscape}
\usepackage{amsmath}
\usepackage{amsfonts}
\usepackage[hidelinks]{hyperref}

\captionsetup{
    font=small,
    labelfont=bf
}

\title{Thesis Structure (6-Chapter Model)\\[0.3cm]\small{Randomized Reward Shaping Deep Reinforcement Learning\\for Self-Adaptive Systems}}
\author{Alireza Kavianifar}
\date{\today}

\begin{document}

\maketitle

\tableofcontents
\newpage

% ============================================================================
% CHAPTER 1: Introduction
% ============================================================================
\vspace*{1cm}
{\centering \huge \textbf{Chapter 1: Introduction} \par}
\addcontentsline{toc}{section}{\textbf{Chapter 1: Introduction}}
\vspace{1cm}

\setcounter{section}{0}
\renewcommand{\thesection}{1.\arabic{section}}
\renewcommand{\thesubsection}{1.\arabic{section}.\arabic{subsection}}

\section{Background and Context: Self-Adaptive Systems}
\subsection{Motivation and Examples}
\subsection{The MAPE-K Feedback Loop}
\subsection{Limitations of Static and Rule-Based Adaptation}

\section{The Challenge of Runtime Uncertainty}
\subsection{Foundations of Uncertainty in Self-Adaptive Systems}
\subsection{Motivating Case Study: The DeltaIoT System (A First Look)}
\subsection{Sources and Impact of Uncertainty in DeltaIoT}

\section{Learning-Based Adaptation: The Promise and Pitfalls of Reinforcement Learning}
\subsection{Why RL for Runtime Decision-Making?}
\subsection{Limits of Classic RL in Large, Discrete Adaptation Spaces}

\section{Problem Statement and Research Questions}
\section{Proposed Solution: Reward-Shaped Deep Reinforcement Learning (RS-DRL)}
\textit{Integrated within the MAPE-K (Monitor-Analyze-Plan-Execute-Knowledge) feedback loop architecture.}

\subsection{Core Idea: Randomized Reward Reshaping ($\rho$-factor) during Experience Replay}
\textit{Controlled by a reshaping factor $\rho$ determining the proportion of failed experiences to reshape.}

\section{Key Contributions of the Thesis}
\section{Thesis Outline}

\newpage

% ============================================================================
% CHAPTER 2: Foundations and State of the Art
% ============================================================================
\vspace*{1cm}
{\centering \huge \textbf{Chapter 2: Foundations and State of the Art} \par}
\addcontentsline{toc}{section}{\textbf{Chapter 2: Foundations and State of the Art}}
\vspace{1cm}

\setcounter{section}{0}
\renewcommand{\thesection}{2.\arabic{section}}
\renewcommand{\thesubsection}{2.\arabic{section}.\arabic{subsection}}

\section{Foundations of Self-Adaptive Systems}
\subsection{Core Concepts and Architectural Patterns}
\subsection{Modeling and Managing Uncertainty in SAS}

\section{Reinforcement Learning for System Adaptation}
\subsection{Fundamentals of RL (MDPs, Q-Learning)}
\subsection{Deep Reinforcement Learning for Large Discrete Action Spaces (DQN)}
\subsection{The Concept of ``Continuous'' Learning in this Thesis (Streaming Adaptation Cycles)}

\section{Related Work: Techniques for Managing Large Adaptation Spaces}
\subsection{Machine Learning-Based Methods (e.g., ML4EAS)}
\subsection{Search-Based Methods (e.g., Genetic Algorithms)}
\subsection{Reinforcement Learning-Based Methods (e.g., Epsilon-Greedy DQN, DLASER)}
\subsection{Reward Shaping and Experience Relabeling (HER, Soft HER)}

\section{Synthesis and Identification of the Research Gap}

\newpage

% ============================================================================
% CHAPTER 3: System Model and Problem Formalization
% ============================================================================
\vspace*{1cm}
{\centering \huge \textbf{Chapter 3: System Model and Problem Formalization} \par}
\addcontentsline{toc}{section}{\textbf{Chapter 3: System Model and Problem Formalization}}
\vspace{1cm}

\setcounter{section}{0}
\renewcommand{\thesection}{3.\arabic{section}}
\renewcommand{\thesubsection}{3.\arabic{section}.\arabic{subsection}}

\section{Motivating Case Study: The DeltaIoT System}
\subsection{System Overview and Sources of Uncertainty}
\subsection{Adaptation Space and Configuration Representation}

\section{A Formal Model for Self-Adaptive Decision Making under Uncertainty}
\subsection{State Space Definition ($S = \{Energy, Packet Loss, Latency\}$)}
\subsection{Action Space: Discrete Adaptation Options}
\subsection{Reward Function Formulation for Single and Multi-Objective Goals}
\subsubsection{Equation 1: Multi-Objective Reward (Article Eq. 11)}
\subsubsection{Equation 2: Single-Objective Reward (Article Eq. 12)}
\subsubsection{Table 1: Reward Formulation for TTS and TT Goals (Article Table 2)}

\section{Problem Statement: Learning an Optimal Adaptation Policy with Continuous Runtime Operation}

\newpage

% ============================================================================
% CHAPTER 4: The RS-DRL Method: Architecture and Algorithm
% ============================================================================
\vspace*{1cm}
{\centering \huge \textbf{Chapter 4: The RS-DRL Method: Architecture and Algorithm} \par}
\addcontentsline{toc}{section}{\textbf{Chapter 4: The RS-DRL Method: Architecture and Algorithm}}
\vspace{1cm}

\setcounter{section}{0}
\renewcommand{\thesection}{4.\arabic{section}}
\renewcommand{\thesubsection}{4.\arabic{section}.\arabic{subsection}}

\section{Overview of the Proposed Approach}

\section{System Architecture: Integrating RS-DRL with the MAPE-K Loop}
\subsection{Managed System and Managing System}
\subsection{Detailed Interaction of the 9-Stage Adaptation Loop}

\section{The RS-DRL Algorithm}
\subsection{Core Algorithm: DQN for Streaming Adaptation Cycles}
\textit{Addressing the limits of classic RL through continuous, streaming learning interaction.}
\subsection{The Reward Shaping Function (Algorithm 2)}
\subsection{Comparison to Hindsight Experience Replay (HER) and other baselines}

\section{Key Implementation Details for Reproducibility}
\subsection{State Construction from Monitored Data}
\subsection{Runtime Verification with UPPAAL-SMC and ActivFORMS}
\subsection{Hyperparameter Selection and Optimization (Grid Search and rho Tuning)}

\newpage

% ============================================================================
% CHAPTER 5: Experimental Evaluation
% ============================================================================
\vspace*{1cm}
{\centering \huge \textbf{Chapter 5: Experimental Evaluation} \par}
\addcontentsline{toc}{section}{\textbf{Chapter 5: Experimental Evaluation}}
\vspace{1cm}

\setcounter{section}{0}
\renewcommand{\thesection}{5.\arabic{section}}
\renewcommand{\thesubsection}{5.\arabic{section}.\arabic{subsection}}

\section{Evaluation Goals and Research Questions}

\section{Experimental Setup}
\subsection{Case Study Instances: DeltaIoTv1, v1.1, v2}
\subsection{Datasets: Environmental Traces and Drift Parameters}
\subsection{Implementation Platform: Gymnasium, Stable-Baselines3}

\section{Baselines and Comparative Methods}
\subsection{Epsilon-Greedy DQN}
\subsection{RS-DRL with HER / Soft HER}
\subsection{DLASER, DLASER+, and ML4EAS}
\subsection{Exhaustive Search Reference Baseline}
\subsection{Justification for Baseline Selection: Limitations and Comparison Protocols}

\section{Evaluation Metrics}
\subsection{Learning Performance: Asymptotic Performance, Time to Threshold, and Total Performance}
\subsection{System Performance: Packet Loss, Latency, Energy Consumption}

\section{Results and Analysis}
\subsection{Learning Curves and Convergence Behavior}
\subsection{Per-Objective Performance Comparison}
\subsection{Comparative Evaluation Against All Baselines}
\subsection{Ablation Study: Sensitivity to the Reshaping Factor (rho)}
\subsection{Scalability Analysis (216 $\rightarrow$ 4,096 adaptation options)}

\section{Statistical Validation (t-tests and p-values)}

\newpage

% ============================================================================
% CHAPTER 6: Discussion, Future Work, and Conclusions
% ============================================================================
\vspace*{1cm}
{\centering \huge \textbf{Chapter 6: Discussion, Future Work, and Conclusions} \par}
\addcontentsline{toc}{section}{\textbf{Chapter 6: Discussion, Future Work, and Conclusions}}
\vspace{1cm}

\setcounter{section}{0}
\renewcommand{\thesection}{6.\arabic{section}}
\renewcommand{\thesubsection}{6.\arabic{section}.\arabic{subsection}}

\section{Discussion of Key Findings}
\subsection{Theorizing the Mechanism: Why RS-DRL Works}
\subsection{Analysis of Trade-offs (e.g., Latency vs. Packet Loss)}
\subsection{Implications for Self-Adaptive Systems Engineering}

\section{Limitations and Threats to Validity}
\subsection{Internal, External, and Construct Validity}
\subsection{Limitations of the RS-DRL Approach (e.g., dependence on rho, discrete action spaces)}

\section{Directions for Future Research}
\subsection{Towards Adaptive Reshaping Strategies (Self-Tuning rho)}
\subsection{Extending RS-DRL to Continuous and Safety-Critical Domains}
\subsection{Cross-Domain Validation and Theoretical Analysis}

\section{Conclusion}
\subsection{Summary of the Research Problem and Proposed Solution}
\subsection{Recap of Contributions and Empirical Evidence}
\subsection{Final Remarks}

\newpage

% ============================================================================
% REFERENCES & APPENDICES
% ============================================================================
\section*{References}
\addcontentsline{toc}{section}{References}

\section*{Appendix A --- Full Process Workflows}
\addcontentsline{toc}{section}{Appendix A --- Full Process Workflows}

\section*{Appendix B --- Full Algorithm Listings}
\addcontentsline{toc}{section}{Appendix B --- Full Algorithm Listings}

\section*{Appendix C --- Extended Experimental Results}
\addcontentsline{toc}{section}{Appendix C --- Extended Experimental Results}

\section*{Appendix D --- Implementation Artifacts \& Replication Package}
\addcontentsline{toc}{section}{Appendix D --- Implementation Artifacts \& Replication Package}

\end{document}
